% TikZ 图：电磁感应变式训练读图提示
% Day 23：电磁感应
\begin{tikzpicture}[scale=0.78,>=Latex,line cap=round,line join=round]
  \path[use as bounding box] (-4.20,-1.35) rectangle (4.20,1.22);

  \begin{scope}[shift={(-2.35,0)}]
    \node[font=\small] at (0,1.02) {第 3 题};
    \draw[thick] (-1.15,0.48) -- (1.10,0.48);
    \draw[thick] (-1.15,-0.48) -- (1.10,-0.48);
    \draw[very thick,red] (-0.18,-0.48) -- (-0.18,0.48);
    \draw[->,red,thick] (-0.18,0) -- (0.55,0) node[right] {$v$};
    \foreach \x in {-0.90,-0.45,0,0.45,0.90} {
      \foreach \y in {-0.24,0.24} {
        \draw[blue] (\x-0.04,\y-0.04) -- (\x+0.04,\y+0.04);
        \draw[blue] (\x-0.04,\y+0.04) -- (\x+0.04,\y-0.04);
      }
    }
    \node[font=\scriptsize] at (0,-1.05) {导体棒切割磁感线};
  \end{scope}

  \begin{scope}[shift={(2.35,0)}]
    \node[font=\small] at (0,1.02) {第 4 题};
    \draw[thick,fill=blue!4] (-1.15,-0.58) rectangle (1.15,0.58);
    \foreach \x in {-0.72,-0.24,0.24,0.72} {
      \foreach \y in {-0.25,0.25} {
        \draw[blue] (\x-0.04,\y-0.04) -- (\x+0.04,\y+0.04);
        \draw[blue] (\x-0.04,\y+0.04) -- (\x+0.04,\y-0.04);
      }
    }
    \draw[thick,orange] (-0.34,-0.27) rectangle (0.34,0.27);
    \draw[->,orange,thick] (-0.05,0) -- (0.75,0) node[right] {$v$};
    \node[font=\scriptsize] at (0,-1.05) {磁通量不变};
  \end{scope}
\end{tikzpicture}
